\begin{textsample}{POS Dim 2 – human – Score 10.00 – rj/rj\_ef\_matematica\_1\_p1.txt}  \label{ex:f2_pos_002}
A Matemática é um importante componente na construção da cidadania, \textbf{recursos} tecnológicos e nos conhecimentos científicos.
Ela fornece \textbf{instrumentos} eficazes para compreender e atuar no mundo que nos cerca, é uma ferramenta essencial na solução de vários tipos de problemas, de vários ramos da ciência.
Suas aplicações são inúmeras, tais como, fenômenos naturais, ciências sociais, ciências biológicas, entre outras.
Os argumentos utilizados para amparar a existência do ensino da Matemática escolar são vários.
Em toda Educação Básica, essa área tem sido considerada \textbf{capaz} de contribuir enormemente na formação intelectual do estudante, seja por sua grande aplicação na \textbf{sociedade} contemporânea, seja pelas suas potencialidades na formação de cidadãos críticos, cientes de suas responsabilidades sociais.
A matemática concebe a essência do que é chamado de pensamento moderno e que a \textbf{partir} do século XVII se estendeu por todo mundo com crescente importância.
De acordo com D’Ambrosio ( 1993 ) é a única disciplina escolar que é ensinada praticamente da mesma forma e com o mesmo conteúdo para todos os estudantes do mundo.

% matched lemmas: capaz, instrumento, partir, recurso, sociedade
\end{textsample}
